%%% Hlavní soubor. Zde se definují základní parametry a odkazuje se na ostatní části. %%%

%% Verze pro jednostranný tisk:
% Okraje: levý 40mm, pravý 25mm, horní a dolní 25mm
% (ale pozor, LaTeX si sám přidává 1in)
\documentclass[12pt,a4paper]{report}
\setlength\textwidth{145mm}
\setlength\textheight{247mm}
\setlength\oddsidemargin{15mm}
\setlength\evensidemargin{15mm}
\setlength\topmargin{0mm}
\setlength\headsep{0mm}
\setlength\headheight{0mm}
% \openright zařídí, aby následující text začínal na pravé straně knihy
\let\openright=\clearpage

\usepackage[utf8]{inputenc}

%% Ostatní balíčky
\usepackage{graphicx}
\usepackage{amsthm}

\usepackage{upgreek}

\usepackage{listings}
\usepackage{tikz}
\usepackage{xyling}

\usepackage[center]{caption}

\lstdefinelanguage{CSharp}[Sharp]{C}
{morekeywords={from,where,join,on,equals,orderby,descending,group,by,let,select},columns=flexible,showstringspaces=false}

\lstset{language=CSharp}

%% Balíček hyperref, kterým jdou vyrábět klikací odkazy v PDF,
%% ale hlavně ho používáme k uložení metadat do PDF (včetně obsahu).
%% POZOR, nezapomeňte vyplnit jméno práce a autora.
\usepackage[unicode]{hyperref}   % Musí být za všemi ostatními balíčky
\hypersetup{pdftitle=.NET library for the MediaWiki API}
\hypersetup{pdfauthor=Petr Onderka}

\usepackage[all]{hypcap}

%%% Drobné úpravy stylu

% Tato makra přesvědčují mírně ošklivým trikem LaTeX, aby hlavičky kapitol
% sázel příčetněji a nevynechával nad nimi spoustu místa. Směle ignorujte.
\makeatletter
\def\@makechapterhead#1{
  {\parindent \z@ \raggedright \normalfont
   \Huge\bfseries \thechapter. #1
   \par\nobreak
   \vskip 20\p@
}}
\def\@makeschapterhead#1{
  {\parindent \z@ \raggedright \normalfont
   \Huge\bfseries #1
   \par\nobreak
   \vskip 20\p@
}}
\makeatother

% Toto makro definuje kapitolu, která není očíslovaná, ale je uvedena v obsahu.
\def\chapwithtoc#1{
\chapter*{#1}
\addcontentsline{toc}{chapter}{#1}
}

\begin{document}

% Trochu volnější nastavení dělení slov, než je default.
\lefthyphenmin=2
\righthyphenmin=2

%%% Titulní strana práce

\pagestyle{empty}
\begin{center}

\large

Charles University in Prague

\medskip

Faculty of Mathematics and Physics

\vfill

{\bf\Large BACHELOR THESIS}

\vfill

\centerline{\mbox{\includegraphics[width=60mm]{img/logo.eps}}}

\vfill
\vspace{5mm}

{\LARGE Petr Onderka}

\vspace{15mm}

% Název práce přesně podle zadání
{\LARGE\bfseries .NET library for the MediaWiki API}

\vfill

% Název katedry nebo ústavu, kde byla práce oficiálně zadána
% (dle Organizační struktury MFF UK)
Department of Theoretical Computer Science and Mathematical Logic

\vfill

\begin{tabular}{rl}

Supervisor of the bachelor thesis: & Tomáš Petříček \\
\noalign{\vspace{2mm}}
Study programme: & Computer Science \\
\noalign{\vspace{2mm}}
Specialization: & General Computer Science \\
\end{tabular}

\vfill

% Zde doplňte rok
Prague 2012

\end{center}

\newpage

%%% Následuje vevázaný list -- kopie podepsaného "Zadání bakalářské práce".
%%% Toto zadání NENÍ součástí elektronické verze práce, nescanovat.

%%% Na tomto místě mohou být napsána případná poděkování (vedoucímu práce,
%%% konzultantovi, tomu, kdo zapůjčil software, literaturu apod.)

\openright

\noindent
Dedication.

\newpage

%%% Strana s čestným prohlášením k bakalářské práci

\vglue 0pt plus 1fill

\noindent
I declare that I carried out this bachelor thesis independently, and only with the cited
sources, literature and other professional sources.

\medskip\noindent
I understand that my work relates to the rights and obligations under the Act No.
121/2000 Coll., the Copyright Act, as amended, in particular the fact that the Charles
University in Prague has the right to conclude a license agreement on the use of this
work as a school work pursuant to Section 60 paragraph 1 of the Copyright Act.

\vspace{10mm}

\hbox{\hbox to 0.5\hsize{%
In ........ date ............
\hss}\hbox to 0.5\hsize{%
signature of the author
\hss}}

\vspace{20mm}
\newpage

%%% Povinná informační strana bakalářské práce

\vbox to 0.5\vsize{
\setlength\parindent{0mm}
\setlength\parskip{5mm}

Název práce:
.NET knihovna pro MediaWiki API
% přesně dle zadání

Autor:
Petr Onderka

Katedra:  % Případně Ústav:
Katedra teoretické informatiky a matematické logiky 
% dle Organizační struktury MFF UK

Vedoucí bakalářské práce:
Mgr. Tomáš Petříček, University of Cambridge
% dle Organizační struktury MFF UK, případně plný název pracoviště mimo MFF UK

Abstrakt:
% abstrakt v rozsahu 80-200 slov; nejedná se však o opis zadání bakalářské práce

Klíčová slova:
% 3 až 5 klíčových slov

\vss}\nobreak\vbox to 0.49\vsize{
\setlength\parindent{0mm}
\setlength\parskip{5mm}

Title:
.NET library for the MediaWiki API
% přesný překlad názvu práce v angličtině

Author:
Petr Onderka

Department:
Department of Theoretical Computer Science and Mathematical Logic
% dle Organizační struktury MFF UK v angličtině

Supervisor:
Mgr. Tomáš Petříček, University of Cambridge
% dle Organizační struktury MFF UK, případně plný název pracoviště
% mimo MFF UK v angličtině

Abstract:
% abstrakt v rozsahu 80-200 slov v angličtině; nejedná se však o překlad
% zadání bakalářské práce

Keywords:
% 3 až 5 klíčových slov v angličtině

\vss}

\newpage

%%% Strana s automaticky generovaným obsahem bakalářské práce. U matematických
%%% prací je přípustné, aby seznam tabulek a zkratek, existují-li, byl umístěn
%%% na začátku práce, místo na jejím konci.

\openright
\pagestyle{plain}
\setcounter{page}{1}
\tableofcontents

%%% Jednotlivé kapitoly práce jsou pro přehlednost uloženy v samostatných souborech
\chapter*{Introduction}
\addcontentsline{toc}{chapter}{Introduction}

\chapter{The goal}

The goal of the LinqToWiki library is to be able to express requests to the MediaWiki API
in a way that is readable, discoverable and checked by the compiler for correctness as much as possible.

This is achieved by generating classes specific for each module and using them in LINQ queries.

The two most commonly used variants of LINQ are LINQ to Objects, and various versions of LINQ for SQL databases.
LINQ to Objects is used for querying in-memory data, like arrays.
There are several widely-used libraries for accessing SQL databases using LINQ, including LINQ to SQL, LINQ to Entities and NHibernate.

In all versions of LINQ, the queries look the same. For example:

\begin{lstlisting}
from product in allProducts
where product.Price > 500
   && product.InStock
join category in categories on product.Category equals category
orderby product.Price
select product.Name
\end{lstlisting}

A query like this is translated into a sequence of method calls that take their parameters in the form of lambda expressions.
For example, the \lstinline{where} part of the above query is translated into:

\begin{lstlisting}
Where(product => product.Price > 500 && product.InStock)
\end{lstlisting}

The commonalities between LINQ to Objects and SQL LINQ libraries are that the full range of operators are available
and that all properties are available in all of them.

The situation with the MediaWiki API is different in several ways:

\begin{enumerate}
\item It doesn't support queries represented by many of the LINQ operators, including \lstinline{join} and \lstinline{group by}.
\item Some of the modules don't support sorting, some do. Of those that do support sorting, some allow specifying the sort key, others only direction.
\item The sets of properties that are available for filtering, sorting and selecting are all different.
\item There are modules used for queries about a set of pages. Those pages can be from a hard-coded list or a result from some other module.
\item There are also parameters that don't fit into the LINQ model well. Some of them are required, some are not.
\end{enumerate}

The goal is to be able to represent all valid queries, while invalid queries should cause a compile-time error.

Specifically, unsupported operators (like \lstinline{join} and \lstinline{group by}) should cause an error for all modules,
while the \lstinline{orderby} clause should cause an error only for the modules that don't support sorting.

Also, all operators should support only those properties that are actually supported by the API.
So, for example for the \verb,blocks, module, the following query should compile and execute fine:

\begin{lstlisting}
from block in wiki.Blocks()
where block.Ip == "8.8.8.8"
orderby block descending
select block.ById
\end{lstlisting}

While the following query should cause three errors:

\begin{lstlisting}
from block in wiki.Blocks()
where block.ById == 1234
orderby block.Expiry descending
select block.Ip
\end{lstlisting}

This is because limiting the query by the blocked IP address, sorting without specifying the key and selecting the ID of the user who performed the block are all allowed.
On the other hand, limiting by the ID of the user who performed the block, sorting by the expiration date and selecting the IP address are all impossible.
(Actually selecting the IP address of the blocked user is possible, but the information is contained in properties with different names.)
\chapter{MediaWiki API}

MediaWiki is an open source wiki system.
It is written in the PHP programming language and uses a relational database to store its data, usually MySQL.
It is maintained by the WikiMedia Foundation, who also runs some of the biggest wikis, including Wikipedia and Wiktionary.
It is also used by many others, including Wikia, who runs many small wikis for various interests and
the unofficial wiki of the Faculty of Mathematics and Physics, $\upomega\upiota\upkappa\upiota$.matfyz.cz.

There are several ways to programmatically access the database of some MediaWiki wiki.
First, it's possible to directly access the database using SQL.
This usually requires access to the server that runs the database, so it's not available in many cases.
For WikiMedia wikis, a read-only access to most data, but excluding article texts,
is available for registered users of Toolserver, run by WikiMedia Deutschland.

Mostly specific to WikiMedia wikis is also another option: database dumps. These are files that contain
dumps of some tables of the wikis. Their disadvantages are that the newest dump is usually several days or weeks old
and that the files can be huge, which is impractical for getting information about a small number of pages.

Last, but not least, is the MediaWiki API.
It can be used to remotely access any MediaWiki wiki (unless the API is disabled in the configuration)
using the HTTP protocol.

Parameters for the API request are given in the query string of a GET request or in the body of a POST request
(some modules require the use of POST).
The body of the POST request is usually formatted as \texttt{application/x-www-form-urlencoded},
but file uploads require the use of \texttt{multipart/form-data}.

Some parameters can accept multiple values at once.
In these cases, the values are separated by a pipe character (\texttt{|}).

The response can be in one of the several available formats, the most widely used ones are XML and JSON.

The API is divided into modules and there are two kinds of modules:
``normal'' modules (called ``non-query modules'' in this work) and query modules.

Non-query modules are usually used to perform some action.
For example the \texttt{edit} module can be used to edit a page
and the \texttt{block} module can be used to block a user (it can be used only by a user with sufficient privileges).

Query modules are used for retrieving information about the wiki. There are three types of query modules:

\begin{itemize}
\item \texttt{list} modules: Return contents of various lists.
For example the \texttt{allcategories} module can be used to list all categories on a wiki,
while the \texttt{categorymembers} module can be used to list members of a certain category.
\item \texttt{prop} modules: Return information about a set of pages.
For example, the \texttt{categories} module can be used to retrieve the categories for each page in a given set.
\item \texttt{meta} modules: Return meta information that are not directly associated with pages.
For example the \texttt{userinfo} module can be used to retrieve information about the currently logged-in user.
\end{itemize}

For \texttt{prop} modules, the set of pages they operate on can be specified directly using page titles or page IDs.

Another option is to use some other module (usually a \texttt{list} module) as a so called ``generator''.
This way, one can for example retrieve the categories of all pages in a specific category,
by using the \texttt{categorymembers} module as a generator for the \texttt{categories} module.

The API is also extensible: MediaWiki extensions can add their own modules and modify some behavior of the existing modules.

An example of an API request URI and a response in the XML format is in Figure \ref{API example}.

\begin{figure}[htbp]
\texttt{http://en.wikipedia.org/w/api.php}~\texttt{?}\
\texttt{format}~\texttt{=}~\texttt{xml}~\texttt{\&}
\texttt{action}~\texttt{=}~\texttt{query}~\texttt{\&}
\texttt{list}~\texttt{=}~\texttt{categorymembers}~\texttt{\&}
\texttt{cmtitle}~\texttt{=}~\texttt{Category:Query\%20languages}~\texttt{\&}
\texttt{cmprop}~\texttt{=}~\texttt{title}~\texttt{\&}
\texttt{cmtype}~\texttt{=}~\texttt{page}~\texttt{\&}
\texttt{cmdir}~\texttt{=}~\texttt{descending}~\texttt{\&}
\texttt{cmlimit}~\texttt{=}~\texttt{5}

\begin{lstlisting}[language=xml]
<?xml version="1.0"?>
<api>
  <query>
    <categorymembers>
      <cm ns="0" title="YQL (programming language)" />
      <cm ns="0" title="Yahoo! query language" />
      <cm ns="0" title="XQuery" />
      <cm ns="0" title="XPath" />
      <cm ns="0" title="XBase++" />
    </categorymembers>
  </query>
  <query-continue>
    <categorymembers cmcontinue="page|5842415345|572327" />
  </query-continue>
</api>
\end{lstlisting}

\caption{An example of an API request and a response}
\label{API example}
\end{figure}

\section{Paging}

Because the results of the API queries can contain thousands and sometimes even millions of entries,
the responses are limited.
For most modules, the default limit (when it is not specified as a parameter) is ten entries per page
and the default maximum is 500~entries for normal users.
For users with the \texttt{apihighlimits} right, the limits are raised, usually to 5000~entries per page.

In the \texttt{limit} parameter, one can specify either the exact value,
or the special value \texttt{max}, which means the maximum allowed for the current user.

To get the data from the following page, one has to use a value specified in the \texttt{query-continue}
element in the result (see Figure \ref{API example} again).
The value in this element is a transparent identifier of the next page.

The advantage of this system when compared with the conventional paging systems of numbering pages
or using numeric offsets is that it avoids missing entries and duplicates when the result
changes while retrieving the pages.

The API has no notion of transactions, so it's not possible to get fully consistent results
that would correspond to an exact moment in time.

But thanks to this paging system, one can be certain that an entry that should be in the result set
during retrieving of all of the pages will actually be present in the result set exactly once.

The situation gets even more complicated when using a \texttt{prop} module with another modules as a generator.
That is because both modules have their own paging.

When such a request is made, the first response will contain a limited number of items from the generator
and a limited number of results from the \texttt{prop} module for those items.
To retrieve the next set of items from the generator, one has to use the \texttt{query-continue} for the generator
(called ``primary paging'' in this work).
To retrieve the next set of results for the items from the first result,
one has to use the \texttt{query-continue} for the \texttt{prop} module (called ``secondary paging'' here).

For an example, see Figure \ref{paging}.
It shows how the paging might work when using the \texttt{allpages} module as a generator,
together with the \texttt{prop} module \texttt{categories}.
The \texttt{query-continue} elements are not shown in the figure.

\begin{figure}[htbp]
\begin{center}
\begin{tikzpicture}
\path (0,12) node(11) {
\begin{minipage}{150pt}
\begin{lstlisting}[language=xml,frame=single]
<query>
  <pages>
    <page title="A" />
    <page title="B">
      <categories>
        <cl title="X" />
        <cl title="Y" />
      </categories>
    </page>
  </pages>
</query>
\end{lstlisting}
\end{minipage}
} (7,12) node(12) {
\begin{minipage}{150pt}
\begin{lstlisting}[language=xml,frame=single]
<query>
  <pages>
    <page title="A">
      <categories>
        <cl title="X" />
        <cl title="Y" />
      </categories>
    </page>
    <page title="B" />
  </pages>
</query>
\end{lstlisting}
\end{minipage}
} (11,12) node(13) {\dots};

\path (0,4) node(21) {
\begin{minipage}{150pt}
\begin{lstlisting}[language=xml,frame=single]
<query>
  <pages>
    <page title="C">
        <cl title="X" />
        <cl title="Y" />
    </page>
    <page title="D" />
  </pages>
</query>
\end{lstlisting}
\end{minipage}
} (7,4) node(22) {
\begin{minipage}{150pt}
\begin{lstlisting}[language=xml,frame=single]
<query>
  <pages>
    <page title="C">
      <categories>
        <cl title="Z" />
      </categories>
    </page>
    <page title="D">
      <categories>
        <cl title="X" />
      </categories>
    </page>
  </pages>
</query>
\end{lstlisting}
\end{minipage}
};

\path (0,0) node(31) {\dots};

\draw[->] (11) -- (12);
\draw[->] (12) -- (13);
\draw[->] (11) -- (21);
\draw[->] (21) -- (22);
\draw[->] (21) -- (31);

%TODO: divné mezery
\end{tikzpicture}
\end{center}

\caption{An example of primary and secondary paging}
\label{paging}

\end{figure}

\section{The \texttt{paraminfo} module}

\chapter{LINQ and expression trees}

LINQ, short for Language Integrated Query, is a feature of the C\# programming language
and the .Net framework that can be used for querying of various data sources.
It uses higher-order functions and lambda expressions to achieve a readable declarative syntax.

LINQ consists of a set of so called “standard query operators”:
methods that are used to perform the query operations on a given source.
Also, a special syntax (called “query expressions”), similar to SQL queries, is available
for some of those operators.
The compiler translates a query expression into a set of calls to standard query operators,
using lambda expressions and anonymous types.

Anonymous types are types that don't have to be explicitly declared;
they are used in similar situations as tuples in functional programming.
An instance of an anonymous type is created by using the \lstinline{new} keyword
without specifying the type of the object to create.

For example the following query expression (as seen in Chapter \ref{goal}):

\begin{lstlisting}
from product in allProducts
where product.Price > 500
   && product.InStock
join category in categories on product.Category equals category
orderby product.Price
select product.Name
\end{lstlisting}

Is translated into the following method calls:

\begin{lstlisting}
allProducts
    .Where(product => product.Price > 500 && product.InStock)
    .Join(
        categories, product => product.Category, category => category,
        (product, category) => new { product, category })
    .OrderBy(t => t.product.Price)
    .Select(t => t.product.Name)
\end{lstlisting}

The parameter~\lstinline,t, is a so called “transparent identifier”, it is used to transfer a set of variables from one method call to another.

The LINQ library also contains methods that do not have a corresponding representation in query expressions. Some examples of those are \lstinline,Aggregate(),, \lstinline,Sum(), and \lstinline,ToList(),.

Many of the basic query operators also correspond to a well-known higher-order function from functional programming. See Figure \ref{LINQ methods} for comparison of some of the LINQ query operators, query expression clauses, and higher-order functions.

\begin{figure}[htbp]
\begin{tabular}{lll}
query operator & query expression clause & higher-order function \\
\lstinline,Select(), & \lstinline,select,, \lstinline,let, & map \\
\lstinline,Where(), & \lstinline,where, & filter \\
\lstinline,SelectMany(), & second and following \lstinline,from, & bind \\
\lstinline,Aggregate(), & & fold \\
\lstinline,Join(), & \lstinline,join, & \\
\lstinline,OrderBy(),,\cr \lstinline,OrderByDescending(), & \lstinline,orderby, & \\
\lstinline,GroupBy(), & \lstinline,group by, & \\
\lstinline,Sum(), \\
\lstinline,First(), \\
\lstinline,ToList(), \\
\end{tabular}

\caption{Comparison between LINQ query operators, query expression clauses, and higher-order functions}
\label{LINQ methods}
\end{figure}

Usually, lambda expressions are compiled into normal methods and passed to the query operator methods as delegates
(which are similar to function pointers in~C or first-class functions in functional languages).
But this would not be suitable for querying of sources that are not in-memory collections,
because the query has to be translated into another form,
like an SQL query or a set of parameters for the MediaWiki API.

Because of this, a lambda expression in C\# can be also compiled into another form:
an expression tree.
Expression tree is an object that represents the given lambda expression
in a form similar to an abstract syntax tree.
This object can be programmatically accessed and manipulated,
which allows translation of LINQ queries into other forms, such as SQL queries.
An expression tree can also be compiled into a delegate.

For an example of an expression tree, see Figure \ref{Expression tree}.

\begin{figure}[htbp]
\begin{center}
\Treek[2]{4}{
 & & & \K{Binary: AndAlso} \B{dll}_{\textnormal{Left}} \B{dr}^{\textnormal{Right}} \\
 & \K{Binary: GreaterThan} \B{dl}_{\textnormal{Left}} \B{dr}^{\textnormal{Right}} & & & \K{Member: InStock} \B{d}^{\textnormal{Expression}} \\
 \K{Member: Price} \B{d}_{\textnormal{Expression}} & & \K{Constant: 500} & & \K{Parameter: product} \\
\K{Parameter: product}
}
\end{center}

\caption{The body of the expression tree for the lambda expression \lstinline,product => product.Price > 500 && product.InStock,}
\label{Expression tree}
\end{figure}

The .Net framework contains two implementations of the query expression pattern:
the interfaces \lstinline,IEnumerable<T>, and \lstinline,IQueryable<T>,.
This means that any object that implements one of these two interfaces can be used in a LINQ query.

These two types implement the query expression pattern completely,
so they can be used with any LINQ operator.
Other custom types can implement only part of the query expression pattern,
which would mean only a subset of the LINQ operators are available for such types.

The \lstinline,IEnumerable<T>, interface usually represents an in-memory collection,
so its implementation of the LINQ operators use delegates.
The \lstinline,IQueryable<T>, interface is usually used to represent a remote collection
(such as a table in a relational database),
so its version of the LINQ operators use expression trees.

The \lstinline,IQueryable<T>, interface doesn't perform any translation of expression trees
into the target query language.
What it does is to combine the whole query into one expression tree,
which is then passed to an implementation of \lstinline,IQueryProvider,.

The query provider is then responsible for processing the expression tree
and translating it into its target query language.
If the query is not valid, the query provider will throw an exception at runtime.
\chapter{Roslyn}

Microsoft Roslyn is a new implementation of the C\# compiler written in C\#
(and a VB.NET compiler written in VB.NET).
Its main distinguishing characteristic is that it is “open”:
it can be used as a library to convert between text and a syntax tree,
to manipulate the syntax tree or to interrogate semantic information.

It also integrates itself into the Microsoft Visual Studio IDE
(Integrated Development Environment), where it can be used to perform
custom refactoring actions or to produce custom errors and warnings at compile-time.

Roslyn is currently under development and its latest public version is a CTP
(Community Technology Preview) from October 2011.
In the CTP, the semantic part of the library in completely implemented,
so for example the syntax tree can represent any construct of C\#
and any syntax tree can be translated to and from source code.

On the other hand, the semantic part of the library is not fully implemented in the CTP,
which means that for example some syntax trees won't successfully compile,
even if they represent valid C\# code.

Because of its close relation with Visual Studio,
Roslyn syntax tree is able to represent every feature of C\# with down to character precision,
including “trivia”: parts of code that are not significant for the compiler,
such as whitespace and comments.

Trivia can also be “structured”, that is it can form a small syntax tree of its own.
An example of this structured trivia are XML documentation comments,
that can be used to provide documentation for a piece of code,
which can then be automatically processed.

For an example of a Roslyn syntax tree, see Figure~\ref{Syntax tree}

\begin{figure}[htbp]

\begin{lstlisting}
public abstract CategoryInfoResult CategoryInfo { get; }
\end{lstlisting}

\begin{center}
\Treek[.5]{4}{
& & & & \K{Property}\Below{Declaration} \B[-5]{dllll} \B[-5]{dll} \B[-5]{d} \B[-5]{drr} \B[-5]{drrrr} \\
\K{Public}\Below{Keyword} & & \K{Abstract}\Below{Keyword} & & \K{Identifier}\Below{Name} \B[-5]{d} & & \K{Identifier}\Below{Token:}\Below{CategoryInfo} & & \K{Accessor}\Below{List} \B[-5]{dll} \B[-5]{d} \B[-5]{drr} \\
& & & & \K{Identifier}\Below{Token:}\Below{CategoryInfoResult} & & \K{OpenBrace}\Below{Token} & & \K{GetAccessor}\Below{Declaration} \B[-5]{dl} \B[-5]{dr} & & \K{CloseBrace}\Below{Token} \\
& & & & & & & \K{Get}\Below{Keyword} & & \K{Semicolon}\Below{Token}
}
\end{center}

\caption{An example of a piece of C\# code and its Roslyn syntax tree \\ (trivia not shown)}
\label{Syntax tree}
\end{figure}

Roslyn syntax trees are immutable and they can be created using factory methods from the \lstinline{Syntax} class.
And while not all elements of the syntax tree have to be specified (like the braces of the accessor list),
creating a syntax tree can be quite cumbersome.

For an example of code to manually create the syntax tree from Figure~\ref{Syntax tree}, see Figure~\ref{Roslyn code}.

\begin{figure}[htbp]

\begin{lstlisting}
Syntax.PropertyDeclaration(
  modifiers:
    Syntax.TokenList(
      Syntax.Token(SyntaxKind.PublicKeyword),
      Syntax.Token(SyntaxKind.AbstractKeyword)),
  type: Syntax.ParseTypeName("CategoryInfoResult"),
  identifier: Syntax.Identifier("CategoryInfo"),
  accessorList:
    Syntax.AccessorList(
      accessors:
        Syntax.List(
          Syntax.AccessorDeclaration(
            SyntaxKind.GetAccessorDeclaration,
            semicolonTokenOpt: Syntax.Token(SyntaxKind.SemicolonToken)))));
\end{lstlisting}

\caption{A sample code to manually create a Roslyn syntax tree}
\label{Roslyn code}
\end{figure}
\chapter{The LinqToWiki library}

\section{The LinqToWiki project}

\section{The LinqToWiki.Codegen project}

\section{The linqtowiki-codegen application}

\chapter{Possible improvements}


\chapter*{Conclusion}
\addcontentsline{toc}{chapter}{Conclusion}


%%% Seznam použité literatury
%%% Seznam použité literatury je zpracován podle platných standardů. Povinnou citační
%%% normou pro bakalářskou práci je ISO 690. Jména časopisů lze uvádět zkrácená, ale jen
%%% v kodifikované podobě. Všechny použité zdroje a prameny musí být řádně citovány.

\def\bibname{Bibliography}
\begin{thebibliography}{12}
\addcontentsline{toc}{chapter}{\bibname}

%\bibitem{lamport94}
%  {\sc Lamport,} Leslie.
%  \emph{\LaTeX: A Document Preparation System}.
%  2. vydání.
%  Massachusetts: Addison Wesley, 1994.
%  ISBN 0-201-52983-1.

\bibitem{mediawiki}
 \emph{MediaWiki}.
 MediaWiki.org.
 \url{http://www.mediawiki.org/wiki/MediaWiki}.

\bibitem{mediawiki-api}
 \emph{MediaWiki API}.
 MediaWiki.org.
 \url{http://www.mediawiki.org/wiki/API}.
 
\bibitem{cs-in-depth}
  {\sc Skeet,} Jon.
  \emph{C\# in Depth}.
  2nd edition.
  Stamford: Manning, 2011.
  Part 3, C\#~3: Revolutionizing how we code.
  ISBN 978-1-935182-47-4.
  
\bibitem{roslyn}
 \emph{Microsoft “Roslyn” CTP}.
 MSDN.
 \url{http://msdn.microsoft.com/en-US/roslyn}.
 
\bibitem{linq-to-sql-functions}
 \emph{Data Types and Functions}.
 LINQ to SQL Reference, MSDN Library.
 \url{http://msdn.microsoft.com/en-us/library/bb386970}.

\bibitem{warren}
 {\sc Warren}, Matt.
 \emph{LINQ: Building an IQueryable Provider -- Part III}.
 The Wayward WebLog.
 \url{http://blogs.msdn.com/b/mattwar/archive/2007/08/01/linq-building-an-iqueryable-provider-part-iii.aspx},
 2 August 2007.
 
\bibitem{reflection-emit}
 \emph{Emitting Dynamic Methods and Assemblies}.
 MSDN Library.
 \url{http://msdn.microsoft.com/en-us/library/8ffc3x75}.

\bibitem{codedom}
 \emph{Dynamic Source Code Generation and Compilation}.
 MSDN Library.
 \url{http://msdn.microsoft.com/en-us/library/650ax5cx}.
 
\bibitem{guidelines-for-names}
 \emph{Guidelines for Names}.
 Design Guidelines for Developing Class Libraries, MSDN Library.
 \url{http://msdn.microsoft.com/en-us/library/ms229002}.
 
\bibitem{wikitools}
 \emph{python-wikitools}.
 Google Code.
 \url{http://code.google.com/p/python-wikitools/}.
 
\bibitem{wikifunctions}
 \emph{WikiFunctions}.
 Wikipedia.
 \url{http://en.wikipedia.org/wiki/Wikipedia:WikiFunctions}.
 
\bibitem{linq-to-wikipedia}
 \emph{Linq to Wikipedia}.
 CodePlex.
 \url{http://linqtowikipedia.codeplex.com/}.
 
\end{thebibliography}

%%% Tabulky v bakalářské práci, existují-li.
\chapwithtoc{List of Tables}

%%% Použité zkratky v bakalářské práci, existují-li, včetně jejich vysvětlení.
\chapwithtoc{List of Abbreviations}

%%% Přílohy k bakalářské práci, existují-li (různé dodatky jako výpisy programů,
%%% diagramy apod.). Každá příloha musí být alespoň jednou odkazována z vlastního
%%% textu práce. Přílohy se číslují.
\chapwithtoc{Attachments}

\openright
\end{document}
